% =========================================================================
% Section II -- Related Work.
% The positioning acknowledges prior routing-aware RL assignment explicitly.
% The distinction claimed here is the offline service-constrained setting,
% measurement of evaluator-induced ranking error, and executor-grounded PPO.
% =========================================================================

\section{Related Work}
\label{sec:related}

\subsection{Learned Scheduling with Transport Abstractions}

L2D~\cite{zhang2020l2d} introduced a graph-based constructive policy for job
shops, using the change in a makespan lower bound as a dense reward. Flexible
job-shop extensions use heterogeneous operation--machine graphs and composite
actions~\cite{song2023fjsp}, and subsequent work develops scalable graph-based
schedulers~\cite{moon2024hgs}. Transport-aware formulations additionally model
AGV movement and operation--machine--vehicle decisions using graph,
transformer, or multi-agent architectures
\cite{yuan2026gin,yang2026cosched,matrans2026}.

The cited transport-aware schedulers principally represent motion with fixed
travel times between resources. This supports inexpensive partial-schedule
evaluation but does not reproduce delays from shared-space reservations or
exclusive pickup and delivery service. We therefore build on established 
graph-based constructive schedulers and focus on a different question: whether 
an idealized evaluator accurately predicts executed makespan and preserves 
schedule rankings under collision-aware execution.

\subsection{Integrated Assignment, Routing, and Service Resources}

Scheduling and material handling have long been recognized as coupled
decisions~\cite{bilge1995time,ulusoy1997ga}. Conflict-aware methods integrate
dispatch and route search for capacitated AGVs~\cite{miyamoto2016local},
container-terminal operations~\cite{zhong2020integrated}, and just-in-time
delivery~\cite{nishida2023jit}. Integrated MAPD likewise uses routed delivery
cost to inform assignment and includes a multi-load variant
\cite{chen2021integrated}. These methods optimize routing-aware solutions per
instance rather than learning the offline constructive policy considered here.

Learning-based integration is also established. Li \emph{et al.}
\cite{li2024hierarchical} train a dynamic task-assignment policy whose action
rewards are evaluated by a lower-level conflict-free path generator. Wang
\emph{et al.}~\cite{wang2026collaborative} jointly train task-allocation and
path-planning agents, while a reinforcement-learning hyper-heuristic combines
task sequencing with conflict-avoidance rules but re-runs search for each
instance~\cite{li2024rlhh}.

Multi-agent path finding computes collision-free space--time trajectories on
shared graphs~\cite{stern2019mapf}; lifelong MAPD combines repeated routing
with online task service~\cite{ma2017lifelong}. Learned components also appear
inside the planner, for example as priority policies for rolling-horizon
prioritized planning~\cite{rlrhpp2026}. Commitment timing is also consequential:
deliberate assignment deferral changes future reservations and deadline
performance~\cite{noda2026deferral}. Nor are service resources absent from the
literature; recent multi-subsystem MAPD work explicitly coordinates
single-dock handover stations with finite buffers~\cite{zang2026buffers}.
These studies establish that assignment, routing, commitment, and service
contention can be modeled jointly. The distinction here is the evaluation
question: we apply an idealized decoder and a fixed deployment executor to a
common set of complete offline schedules, making evaluator-induced value and
ranking errors directly measurable.

\subsection{Learning from Terminal Cost}

Policy-gradient methods can optimize sequence-level objectives from a terminal
return~\cite{williams1992reinforce}, and PPO stabilizes policy updates through a
clipped surrogate objective~\cite{schulman2017ppo}. In neural combinatorial
optimization, rollout baselines reduce variance when solution cost is available
only after construction is complete~\cite{kool2019attention}. These methods
address learning from delayed solution-level feedback, but they do not by
themselves determine whether that feedback reflects deployment execution. In
our method, executor grounding applies to the terminal cost rather than to
intermediate transitions. Both training and inference advance partial
schedules with a calibrated surrogate, while the fixed collision-aware
executor evaluates completed schedules. During training, its makespan supplies
the terminal cost; the surrogate-projected state does not define the objective.

Taken together, prior work establishes learned schedule construction,
routing-aware assignment, and explicit treatment of service resources. The
cited literature does not make evaluator fidelity over a common set of complete
offline schedules its primary object of study. We quantify both value and
ranking errors induced by an idealized transport model and train PPO against
executor-evaluated terminal cost while exposing surrogate-projected congestion
information during schedule construction.
